\chapter{Class Reference}
\label{ch:classes}

This chapter documents the firmware's public classes at the design level:
what each class is responsible for, which collaborators it depends on, its
key invariants, and how it fits the system. It complements --- and does
not duplicate --- the generated API documentation, which carries the exact
method signatures. The HTTP JSON endpoints are documented separately in
Chapter~\ref{ch:api}.

\begin{drnote}[How this chapter grows]
Per the development rules, every public, architecturally-significant class
(the drivers, the application services, and the public domain-core types)
gets a section here, tagged \texttt{\textbackslash label\{cls:ClassName\}},
added in the same change that introduces the class. A tooling check keeps
this chapter in step with the code, so it is always current.
\end{drnote}

The firmware is under active development. As each public class lands, its
section appears below, grouped by layer: domain core, application
services, and hardware drivers.

% ------------------------------------------------------------------
% Domain core (Slice 1)
% ------------------------------------------------------------------

\section{FirmwareVersion}\label{cls:FirmwareVersion}
Strong type wrapping the firmware release identifier (e.g.\ \texttt{0.1.0}).
Used by \texttt{HealthStatus} and the \texttt{/api/health} endpoint so
version strings are never passed as bare \texttt{char*} across module
boundaries. Invariant: non-empty at construction.

\section{HealthStatus}\label{cls:HealthStatus}
Immutable health-check DTO returned by \texttt{GET /api/health}. Built via
\texttt{HealthStatus::ok(FirmwareVersion)} in the pure core; JSON
serialisation lives in \texttt{serializeHealthStatusJson()}. No ESP-IDF
headers.

% ------------------------------------------------------------------
% Network shell (Slice 1)
% ------------------------------------------------------------------

\section{SoftApConfig}\label{cls:SoftApConfig}
Immutable value type holding SoftAP parameters (SSID, channel, station
limit). Factory \texttt{setupDefault()} yields SSID \texttt{DigiRadio-setup}
for first-time provisioning.

\section{SoftApHost}\label{cls:SoftApHost}
RAII wrapper around the ESP-IDF Wi-Fi stack that starts and stops the
configured SoftAP. Imperative shell; no business logic.

\section{SetupWebServer}\label{cls:SetupWebServer}
Minimal HTTP server: gzipped setup UI, \texttt{GET /api/health},
\texttt{POST /api/wifi}, and tuner routes (\texttt{/api/tuner/*}).
JSON parsing and serialisation delegate to the pure core; credentials
persist via \texttt{ISecureStore}; tuner operations via
\texttt{tuner::TunerService}.

\section{NetBootstrap}\label{cls:NetBootstrap}
Owns network resources for setup or STA mode.
\texttt{start(store, tuner)} initialises the platform, joins stored Wi-Fi
when credentials exist, or falls back to the \texttt{DigiRadio-setup}
SoftAP. Must outlive \texttt{app\_main} for the process lifetime.

% ------------------------------------------------------------------
% Domain core + secure store (Slice 2)
% ------------------------------------------------------------------

\section{Secret}\label{cls:Secret}
Opaque wrapper for sensitive strings. Buffer is zeroised on destruction;
no \texttt{operator<<} and no plaintext serialisation. Shell code uses
\texttt{usePlaintext()} transiently for driver APIs only.

\section{WifiSsid}\label{cls:WifiSsid}
Validated Wi-Fi SSID (1--32 bytes). Parsed once from untrusted JSON at
the HTTP boundary before persistence or driver calls.

\section{WifiCredentials}\label{cls:WifiCredentials}
Domain value pairing \texttt{WifiSsid} with a \texttt{Secret} password.
Open networks use an empty password; WPA-PSK requires 8--63 characters.

\section{ISecureStore}\label{cls:ISecureStore}
Abstract persistence boundary for credentials at rest. Slice~2 implements
Wi-Fi credential save/load/clear; station list and user credentials arrive
in later slices. Host tests use fakes; the shell uses \texttt{NvsSecureStore}.

\section{StaClient}\label{cls:StaClient}
RAII STA join helper with an explicit connect timeout. Assumes
\texttt{esp\_wifi\_init()} was already called by \texttt{NetBootstrap}.

\section{NvsSecureStore}\label{cls:NvsSecureStore}
\texttt{ISecureStore} implementation backed by an NVS namespace. Passwords
are stored as NVS strings and never logged. Production should enable NVS
encryption using the reserved \texttt{nvs\_keys} partition.

% ------------------------------------------------------------------
% Hardware drivers (Slice 3)
% ------------------------------------------------------------------

\section{IFirmwareBlobReader}\label{cls:IFirmwareBlobReader}
Abstract streaming reader for firmware blobs embedded in flash. The Si4684
driver uses it to HOST\_LOAD patch and application images in bounded SPI
chunks without allocating the full image on the heap. Host tests use the
same interface over in-memory buffers.

\section{EmbeddedBlobReader}\label{cls:EmbeddedBlobReader}
Non-owning \texttt{IFirmwareBlobReader} over a contiguous byte range (linker
symbols from \texttt{EMBED\_FILES} or test data). Implements offset-based
\texttt{read()} with truncation at end-of-blob.

\section{Si4684EmbeddedImages}\label{cls:Si4684EmbeddedImages}
Owns \texttt{EmbeddedBlobReader} views over \texttt{rom\_patch\_016.bin},
\texttt{dab\_firmware.bin}, and \texttt{fm\_firmware.bin}. Constructed once;
passed to \texttt{Si4684Driver}. \texttt{applicationImage(Si4684Band)} selects
DAB or FM.

\section{Si4684Driver}\label{cls:Si4684Driver}
Full SPI driver for the Si4684 DAB+/FM tuner (Chapter~\ref{ch:si4684}).

\texttt{boot(Si4684Band)} runs AN649 HOST\_LOAD, configures I\textsuperscript{2}S
and band-specific properties, and records \texttt{loadedBand()}. After boot the
driver exposes FM tuning/seek/RSQ/RDS, DAB ensemble tuning, digital service list
fetch, and \texttt{startDabService}. Property access (\texttt{setProperty},
\texttt{setVolume}) and diagnostics (\texttt{getPartInfo}, \texttt{getSysState})
are included. Wrong-band calls return \texttt{Si4684Error::WrongBand}.
Register-level opcodes remain private; see \texttt{Si4684Types.hpp} for status
DTOs.

\section{Adau1701Driver}\label{cls:Adau1701Driver}
RAII I2C driver for the ADAU1701 SigmaDSP. \texttt{boot()} asserts RESET\#,
initialises the shared I2C bus, and replays the SigmaStudio export from
\texttt{Firmware/ADAU1701-Firmware/} on every power-up (no EEPROM self-boot
on DigiRadio).

% ------------------------------------------------------------------
% Application services (Slice 4)
% ------------------------------------------------------------------

\section{ITuner}\label{cls:ITuner}
Abstract tuner boundary in the pure core. Defines boot, band selection,
status readout, FM/DAB tuning, service list, playback, and volume without
ESP-IDF types. \texttt{si4684::Si4684Tuner} implements it on device; host
tests can use fakes.

\section{TunerService}\label{cls:TunerService}
Application service exposing intent-level tuner operations to HTTP and
future UI. Holds a reference to \texttt{core::ITuner}, tracks last tune
target and volume, and maps driver failures to \texttt{core::TunerError}.

\section{Si4684Tuner}\label{cls:Si4684Tuner}
\texttt{ITuner} adapter over \texttt{Si4684Driver}. Translates domain calls
into SPI commands and maps \texttt{Si4684Error} to \texttt{TunerError}.
Constructed once in \texttt{HardwareBootstrap} alongside the driver.
